% Cuestión V. El fragmento de película oculto en secreto1.bmp. Cifras de generado/datos-analisis.tex.

El video está en \texttt{secreto1.bmp}, oculto con LSB4 y cifrado con \texttt{aes256} en modo \texttt{ofb} (Cuestión~IV). La contraseña que lo descifra
es \texttt{metadata} seguida de un retorno de carro (\texttt{0x0d}), probablemente porque la cátedra lo generó con un script con finales de línea de Windows:
\begin{lstlisting}
./stegobmp -extract -p secreto1.bmp -out secreto1 -steg LSB4 -a aes256 -m ofb -pass $'metadata\r'
\end{lstlisting}
Se obtiene un video \texttt{.wmv} de \dato{an-secreto1-wmv-bytes} bytes y \dato{an-secreto1-wmv-duracion}~s: una escena de la serie \emph{Bones}
(temporada 5, episodio 2). Angela y Hodgins examinan una foto que la víctima había enviado por correo y Angela nota que <<este archivo es más grande de lo que debería>>;
Hodgins concluye que <<ocultó datos en el archivo>>.

\begin{itemize}
\item \textbf{Portador:} una fotografía digital (la víctima saludando), enviada como adjunto de un correo electrónico.
\item \textbf{Qué se ocultaba:} datos o código con información secreta de la víctima, que aparecen superpuestos sobre la foto al analizarla.
\end{itemize}

La pista que delata el ocultamiento, el tamaño del archivo, es la misma que delata a \texttt{madrid.bmp} (Cuestión~VI): los métodos LSB no cambian el tamaño del portador.
